
The ML reproducibility crisis is well-documented. Raff [2019] found that approximately 50\% of top-venue ML papers from 1984--2017 fail independent reproduction attempts. Subsequent work has proposed normative interventions: checklists requiring code submission and dataset specification \citep{Pineau2021Improving}, and structured dataset documentation frameworks \citep{Gebru2021Datasheets}. What is absent from the literature is an empirical linkage between \textit{metadata-observable signals} --- signals derivable from public APIs without re-running experiments --- and reproducibility failure outcomes.

The theoretical motivation for such a linkage exists across two independent research threads. D'Amour et al. [2022] identified underspecification as a credibility challenge in modern ML: models optimized on highly concentrated benchmark datasets learn dataset-specific artifacts rather than generalizable patterns, producing results that break under independent replication. Gebru et al. [2021] proposed that absent \texttt{intended\textbackslash \_use} and \texttt{out\textbackslash \_of\textbackslash \_scope\textbackslash \_use} fields in dataset documentation increase the risk of out-of-context dataset application. If these mechanisms operate in practice, then metadata-observable proxies --- HuggingFace Hub field-presence scores for documentation completeness, and Herfindahl-Hirschman Index over OpenML run counts for benchmark concentration --- should predict which papers in Raff's corpus fail independent reproduction. Testing this hypothesis would provide the first empirical bridge between the dataset documentation and ML reproducibility literatures.

The gap in existing work is structural: no study has tested this linkage because it requires querying post-2019 APIs against a pre-2018 corpus. We designed a systematic infrastructure feasibility study to determine whether the required data exists before investing in regression analysis.

Our central finding is that HuggingFace Hub and OpenML exhibit strong domain and temporal selection biases aligned with their founding communities, not with the broader benchmark ecosystem of the pre-2018 ML literature. HF Hub emerged from the NLP/transformer community: IMDB, SST-2, AG News, SQuAD, and GLUE components are present; large-scale DL vision datasets (ImageNet, COCO, Pascal VOC, KITTI, Caltech-101, Caltech-256), speech datasets (LibriSpeech, TIMIT, Wall Street Journal), and graph datasets (Cora, Citeseer) return no valid cards. OpenML was designed for classical tabular ML experiments and architecturally cannot host large image or audio datasets: iris, adult, and covertype are present, as are IMDB, 20 Newsgroups, and MNIST; CIFAR, ImageNet, and other DL benchmarks are absent. This is not a pipeline defect --- 18 of 18 tests pass and the pipeline completes against live APIs in under 90 seconds. The coverage gap is in the platforms themselves.

This paper makes the following contributions:

\textbf{C1 (Empirical):} The first systematic quantification of HuggingFace Hub and OpenML coverage for the Raff 2019 reproducibility corpus, establishing that HF Hub achieves 30\% dataset card coverage (15/50 datasets found; all NLP benchmarks; large-scale DL vision, speech, and graph datasets absent) and OpenML achieves 22\% pre-publication temporal coverage (11/50; predominantly classical tabular datasets plus IMDB, 20 Newsgroups, and MNIST). Both fall below the minimum thresholds required for valid regression analysis ($\geq 50$\% and $\geq 70$\%, respectively).

\textbf{C2 (Infrastructure characterization):} Documentation of the domain-temporal selection bias in both APIs: HF Hub's ecosystem reflects its NLP/transformer origin; OpenML's architecture structurally excludes large image, audio, and graph datasets. Any reproducibility auditing study of pre-2018 ML literature using these APIs will be structurally biased toward NLP and classical tabular ML papers, systematically omitting DL vision and speech papers.

\textbf{C3 (Methodological):} Identification of Papers With Code API as the most appropriate primary data source for this class of reproducibility auditing study, motivated by the empirical characterization of HF Hub and OpenML limitations. This is a theoretical recommendation; the claim has not been empirically verified.

\textbf{C4 (Reusable infrastructure):} A validated, 18-test data acquisition pipeline --- comprising RaffParser, HFCoverageChecker, OpenMLTemporalChecker, MetricsAggregator, and Visualizer modules --- that is immediately reusable for a redesigned study.

The remainder of this paper is organized as follows. Section 2 surveys related work in ML reproducibility auditing, dataset documentation, and benchmark concentration. Section 3 describes the pipeline design and experimental protocol. Section 4 describes the experimental setup. Section 5 presents the coverage measurements and domain analysis. Section 6 discusses implications, limitations, and future directions. Section 7 concludes.

